% ХАВСРАЛТ А: Өгөгдлийн Сангийн Схем ба API Тодорхойлолт
\chapter{ХАВСРАЛТ А: ӨГӨГДЛИЙН САНГИЙН СХЕМ БА API ТОДОРХОЙЛОЛТ}
\label{AppendixA}
\pagecolor{white}

%===============================================================
\section*{А.1 PostgreSQL Өгөгдлийн Сангийн Схем}

\textbf{users хүснэгт:}
\begin{lstlisting}[language=SQL, caption={users хүснэгтийн схем}]
CREATE TABLE users (
    id          UUID PRIMARY KEY DEFAULT gen_random_uuid(),
    restaurant_id UUID REFERENCES restaurants(id),
    email       VARCHAR(255) NOT NULL UNIQUE,
    password_hash VARCHAR(255) NOT NULL,
    role        VARCHAR(50) NOT NULL 
                CHECK (role IN ('admin','manager','waiter','chef','cashier')),
    first_name  VARCHAR(100),
    last_name   VARCHAR(100),
    is_active   BOOLEAN DEFAULT true,
    created_at  TIMESTAMPTZ DEFAULT NOW()
);
\end{lstlisting}

\textbf{table\_areas хүснэгт:}
\begin{lstlisting}[language=SQL, caption={table\_areas хүснэгтийн схем}]
CREATE TABLE table_areas (
    id              UUID PRIMARY KEY DEFAULT gen_random_uuid(),
    restaurant_id   UUID NOT NULL REFERENCES restaurants(id),
    name            VARCHAR(100) NOT NULL,
    capacity        INTEGER DEFAULT 4,
    qr_token        VARCHAR(500),
    qr_generated_at TIMESTAMPTZ,
    status          VARCHAR(20) DEFAULT 'available'
                    CHECK (status IN ('available','occupied','reserved')),
    created_at      TIMESTAMPTZ DEFAULT NOW()
);
\end{lstlisting}

\textbf{menu\_items хүснэгт:}
\begin{lstlisting}[language=SQL, caption={menu\_items хүснэгтийн схем}]
CREATE TABLE menu_items (
    id              UUID PRIMARY KEY DEFAULT gen_random_uuid(),
    category_id     UUID NOT NULL REFERENCES menu_categories(id),
    name            VARCHAR(200) NOT NULL,
    description     TEXT,
    price           DECIMAL(10,2) NOT NULL CHECK (price >= 0),
    image_url       VARCHAR(500),
    allergens       TEXT[],
    is_available    BOOLEAN DEFAULT true,
    preparation_time INTEGER DEFAULT 15,
    display_order   INTEGER DEFAULT 0,
    created_at      TIMESTAMPTZ DEFAULT NOW()
);
\end{lstlisting}

\textbf{orders хүснэгт:}
\begin{lstlisting}[language=SQL, caption={orders хүснэгтийн схем}]
CREATE TABLE orders (
    id              UUID PRIMARY KEY DEFAULT gen_random_uuid(),
    table_id        UUID NOT NULL REFERENCES table_areas(id),
    session_id      UUID NOT NULL,
    status          VARCHAR(20) NOT NULL DEFAULT 'pending'
                    CHECK (status IN 
                        ('pending','confirmed','preparing','ready','served','paid','cancelled')),
    special_requests TEXT,
    subtotal        DECIMAL(10,2),
    total           DECIMAL(10,2),
    created_at      TIMESTAMPTZ DEFAULT NOW(),
    updated_at      TIMESTAMPTZ DEFAULT NOW()
);
\end{lstlisting}

%===============================================================
\section*{А.2 API Endpoint-уудын Дэлгэрэнгүй Тодорхойлолт}

\textbf{POST /api/auth/login}\\
Хүсэлт: \{ email, password \}\\
Хариу 200: \{ token, refreshToken, user: \{id, email, role\} \}\\
Хариу 401: \{ error: "Invalid credentials" \}

\textbf{GET /api/tables/:id/menu?t=\{token\}}\\
Хүсэлт: token параметр (JWT)\\
Хариу 200: \{ table: \{id, name\}, sessionId, menu: [categories + items] \}\\
Хариу 401: \{ error: "Invalid or expired QR token" \}

\textbf{POST /api/orders}\\
Хүсэлт: \{ tableId, sessionId, items: [\{menuItemId, quantity, specialRequest\}], specialRequests \}\\
Хариу 201: \{ orderId, status: "pending", estimatedTime \}\\
Хариу 400: \{ error: "Menu item not available" \}

\textbf{PATCH /api/orders/:id/status}\\
Хүсэлт: \{ status \}, Bearer token шаардлагатай\\
Хариу 200: \{ orderId, status, updatedAt \}\\
Хариу 400: \{ error: "Invalid status transition" \}\\
Хариу 403: \{ error: "Insufficient permissions" \}

%===============================================================
\section*{А.3 WebSocket Событийн Тодорхойлолт}

\textbf{Клиентийн илгээдэг событиуд:}
\begin{itemize}
    \item \texttt{join\_guest\_session} \{sessionId\} — Зочны өрөөнд нэгдэх
    \item \texttt{join\_staff} \{restaurantId, role\} — Ажилтны өрөөнд нэгдэх
    \item \texttt{join\_kitchen} \{restaurantId\} — Кухны өрөөнд нэгдэх
\end{itemize}

\textbf{Серверийн цацдаг событиуд:}
\begin{itemize}
    \item \texttt{new\_order} — Шинэ захиалга (restaurant\_\{id\} өрөөнд)
    \item \texttt{order\_status\_changed} — Статус өөрчлөлт (session + staff өрөөнд)
    \item \texttt{menu\_updated} — Цэс шинэчлэгдсэн (restaurant өрөөнд)
    \item \texttt{table\_status\_changed} — Ширээний статус өөрчлөлт
\end{itemize}
